\documentclass[11pt,a4paper]{article}
\usepackage[utf8]{inputenc}
\usepackage[T1]{fontenc}
\usepackage[margin=0.72in]{geometry}
\usepackage{enumitem}
\usepackage{titlesec}
\usepackage{xcolor}
\usepackage{hyperref}

\definecolor{primarycolor}{RGB}{33,37,41}
\definecolor{secondarycolor}{RGB}{0,92,153}
\definecolor{accentcolor}{RGB}{92,103,112}
\hypersetup{colorlinks=true,urlcolor=secondarycolor}
\titleformat{\section}{\large\bfseries\color{primarycolor}}{}{0em}{}
\titlespacing*{\section}{0pt}{1em}{0.35em}
\setlist[itemize]{leftmargin=1.25em,itemsep=0.15em,topsep=0.2em}
\pagestyle{empty}
\setlength{\parindent}{0pt}
\setlength{\parskip}{0.25em}

\begin{document}
\begin{center}
{\LARGE\bfseries MARTA KOWAL-BRANDT}\\
{\large QA Automation Engineer | DevOps / Platform Engineering}\\
Wroclaw, Poland | marta.kowalbrandt@example.com | +48 555 303 303 | gitlab.com/mkb-example | github.com/mkb-dev
\end{center}

\section*{PROFESSIONAL SUMMARY}
QA automation engineer with 7 years of experience, including the last two focused primarily on CI/CD pipelines, test environments and build tooling. Owns CI pipelines and Kubernetes-based ephemeral test environments for three product teams. Reduced average pipeline time from 38 to 14 minutes and helped lower flaky test rates from 7\% to below 2\%. Developing further hands-on infrastructure experience through a personal k3s and GitOps homelab, including Terraform-provisioned Hetzner infrastructure.

\section*{TECHNICAL SKILLS}
CI/CD and containers: GitLab CI, Jenkins, Docker, Docker Compose\\
Infrastructure and platforms: Kubernetes (namespace-level work on a shared cluster), Helm, k3s, Terraform (personal projects), Argo CD, Ansible (basic)\\
Monitoring and data: Prometheus, Grafana, Loki, PostgreSQL\\
Programming and scripting: Python, Bash, Java (test code), SQL\\
Testing: pytest, Playwright, Selenium, Postman\\
Cloud learning: AWS EC2 and S3 basics (personal learning)\\
Other: Linux, Git

\section*{PROFESSIONAL EXPERIENCE}
\textbf{QA Automation Engineer (Test Infrastructure) | Odrava Health Software, Wroclaw | 2021-Present}
\begin{itemize}
\item Own GitLab CI pipelines for three product teams (approximately 25 developers), covering build, unit, integration and end-to-end test stages, and artifact publishing.
\item Reduced average pipeline time from 38 to 14 minutes by parallelising test jobs, caching dependencies and separating end-to-end tests into their own stage.
\item Built short-lived merge-request test environments with Docker Compose and later Helm on the company's shared Kubernetes cluster; write Helm charts and work with namespace-level access.
\item Maintain a pytest and Playwright end-to-end suite of approximately 900 tests; reduced the flaky test rate from 7\% to below 2\%.
\item Wrote Bash and Python tooling to seed test data into PostgreSQL.
\item Wrote the internal guide to how CI works, which is used in onboarding; mentored two QA interns in 2023.
\end{itemize}

\textbf{QA Engineer | Brightmill Software House, Wroclaw | 2019-2021}
\begin{itemize}
\item Tested web and mobile applications through manual and automated testing for client projects.
\item Introduced a Java and Selenium automated regression suite for the largest client project.
\item Configured Jenkins jobs for nightly regression runs.
\end{itemize}

\textbf{Customer Support Specialist (2nd line) | Brightmill Software House | 2018-2019}
\begin{itemize}
\item Investigated customer issues using logs and databases, and reproduced bugs for developers.
\end{itemize}

\section*{SELECTED INFRASTRUCTURE PROJECTS}
\textbf{Homelab k3s Cluster with GitOps | Personal | 2024-2025}
\begin{itemize}
\item Built a three-node k3s cluster on Raspberry Pi 5 devices; deployed applications from a Git repository using Argo CD.
\item Set up Prometheus, Grafana and Loki monitoring, Traefik ingress and cert-manager for TLS.
\item Provisioned infrastructure for a small Hetzner VPS with Terraform; used Ansible at a basic level.
\end{itemize}

\textbf{Merge-Request Ephemeral Environments | Odrava Health | 2023-2024}
\begin{itemize}
\item Set up short-lived environments for each merge request, with the application, PostgreSQL and seeded data deployed via Helm to a shared Kubernetes cluster from GitLab CI.
\item Configured automatic cleanup after merge or 48 hours of inactivity; QA and product teams use the environments to review features, and end-to-end tests run against them.
\end{itemize}

\textbf{CI Pipeline Speed-up | Odrava Health | 2022}
\begin{itemize}
\item Profiled pipelines across three teams, parallelised test jobs, added dependency caching and split end-to-end tests into a separate stage.
\item Reduced average pipeline time from 38 to 14 minutes and monthly CI runner cost by approximately 30\%.
\end{itemize}

\textbf{Flaky Test Dashboard | Odrava Health | 2023}
\begin{itemize}
\item Collected CI test results in PostgreSQL and built Grafana dashboards showing flaky tests by team.
\item Flaky test rates fell from 7\% to below 2\% after teams began addressing the top offenders.
\end{itemize}

\section*{EDUCATION}
BEng Computer Engineering | Silesian Lowlands University of Technology | 2018

\section*{CERTIFICATIONS AND DEVELOPMENT}
ISTQB Foundation Level | 2019\\
Certified Kubernetes Administrator (CKA): studying; exam booked for January 2027

\section*{LANGUAGES}
Polish (native), English (C1), German (A2)
\end{document}